\subsection{Open-world graph completion and stopping criteria}

A useful neighboring literature comes from knowledge-graph completion. Standard link-prediction benchmarks often treat unobserved triples as negatives even though a real knowledge graph is incomplete. Open-world knowledge-graph work instead permits unseen entities and missing true facts \cite{shi2017}, while later analyses show that common completion metrics can behave misleadingly when unknown triples contain undiscovered positives \cite{yang2022}. Progressive knowledge-graph completion makes verification itself part of the construction loop: candidates are mined, verified and then used to update the graph rather than being treated as ground truth at retrieval time \cite{li2024kg}. These systems solve a different problem from RAKL, but they reinforce two boundaries used here: an absent fact is not automatically false, and graph completion is an ongoing verified process rather than a certificate that the external world has been exhausted.

The provenance literature also reaches further into the fixed-point setting than a simple source graph. Recent semiring work extends provenance analysis to first-order logic with negation and studies how model-checking results depend on both positive and negative atomic information \cite{gradel2024}. This narrows the novelty boundary around the present saturation construction: the contribution is not provenance for logical inference or fixed points in isolation, but the joint certificate tying provenance, open-world route coverage, evidence independence, novelty audits, freshness and repeated marginal flatness to a scientific release decision.

Stopping research introduces a second adjacent distinction. Decision-theoretic screening work models the decision to continue reading as a value-of-information problem rather than assuming that a fixed recall target is always the right objective \cite{fletcher2026}. A complementary 2026 study proposes confidence-based stopping rules that ask whether the screened evidence is already sufficient to preserve the eventual review conclusion, not merely whether a nominal recall threshold has been reached \cite{fletcher2026b}. Epistemic saturation is stricter in another direction: it is not a screening heuristic over one ranked document stream, because derivations, counterexamples, novelty boundaries, mechanism owners and discovery routes can all change the state. The common lesson is nevertheless important: a stopping rule should be tied to the epistemic decision it is meant to support, not to the raw number of documents or graph edges processed.

These adjacent results modify the saturation interpretation in a useful way. Flatness of the paper's knowledge vector is a release condition only after the system has checked that the relevant decision coordinates themselves are stable. If a new source would not change any claim, proof, evidence root, contradiction, novelty boundary, assumption, unresolved fiber or discovery route, then adding it is bibliographic redundancy rather than epistemic growth. If it changes even one of those coordinates, the state is not flat regardless of how high the prior screening recall appeared to be.
